\section{D1. Normalized local densities and a proved near-maximizer selection theorem}
\label{sec:r33-d1}

This section distinguishes a change of variables, a conditional local-Laplace
calculation, and a mechanical theorem.  A local theorem for an additive sum
does not already construct phase labels, Morse--Bott manifolds, or an adaptive
controlled local theorem.  The first two levels are proved here with their
inputs visible.  For policy optimization we replace a restatement of the
desired epi-development by checkable power-drop and Holder estimates.

\subsection{The continuous Jacobian is compulsory}

Let $K_N\in\mathbb Z^{d_Z}$ and $S_N\in\mathbb R^{d_R}$ be unnormalized
sums.  Suppose their one-lattice-site density in the central region is
\[
 p_N(k,s)=N^{-(d_Z+d_R)/2}
 \left[\varphi_\Sigma\left(
 \frac{k-Nm_Z}{\sqrt N},\frac{s-Nm_R}{\sqrt N}\right)+o(1)\right].
\]
The reference measure here is counting measure in $k$ and Lebesgue measure
in $s$.  Define $Y_N=S_N/N$ and keep $K_N$ unnormalized.

\begin{proposition}[Raw-sum to empirical-mean conversion]
\label{prop:r33-d1-jacobian}
The density of $(K_N,Y_N)$ is exactly $N^{d_R}p_N(k,Ny)$.
Consequently its central prefactor is $N^{(d_R-d_Z)/2}$, not
$N^{-d_Z/2}$.  The consistent fluctuation vector is
\[
 \left((K_N-Nm_Z)/\sqrt N,\sqrt N(Y_N-m_R)\right).
\]
\end{proposition}
\begin{proof}
The change of variables $s=Ny$ gives $ds=N^{d_R}dy$.  Substitute this in the
probability integral and collect powers.  In particular a standard Gaussian
sample mean has density $(N/2\pi)^{d_R/2}e^{-N|y|^2/2}$, whose integral is
one.  Omitting the positive power fails already in this exact example.
\end{proof}

\subsection{A conditional fibre/cell Laplace calculation}

For a specified label $j$, suppose an actual local density, not merely an
LDP, has been established in a tubular phase neighbourhood:
\begin{equation}\label{eq:r33-d1-local}
 p_{j,N}(k,y)=e^{-NI_j(k/N,y)}N^{\lambda_j-d_Z/2}
                 [a_j(k/N,y)+o(1)].
\end{equation}
Assume the remainder is uniform relative to a positive continuous amplitude;
the minimizers are $\{k_j\}\times\mathcal M_j$, with
$\mathcal M_j$ a compact smooth $r_j$-dimensional manifold; the quadratic form
in the $d_Z+d_R-r_j$ normal directions is uniformly positive; a common
Gaussian bound dominates the rescaled integrands; and the part outside the
tube is negligible at the claimed polynomial order.  These are explicit
local-density hypotheses, not consequences of a central Edgeworth theorem.
For simplicity $K_N$ has span one in each lattice coordinate; a different
span changes the covolume constant, not the powers below.

\begin{theorem}[Leading powers for an exact fibre and a summed cell]
\label{thm:r33-d1-laplace}
Pinning $K_N=Nk_j+O(1)$ and integrating the continuous tube gives
\[
 P(\mathcal C_{j,N}^{\rm fib})
 =e^{-N\gamma_j}N^{-\kappa_j^{\rm fib}}[C_j^{\rm fib}+o(1)],
 \qquad
 \kappa_j^{\rm fib}=d_Z/2-\lambda_j+(d_R-r_j)/2.
\]
Summing the lattice coordinates over a central window of width
$N^{1/2+\epsilon}$ instead gives
\[
 P(\mathcal C_{j,N}^{\rm cell})
 =e^{-N\gamma_j}N^{-\kappa_j^{\rm cell}}[C_j^{\rm cell}+o(1)],
 \qquad
 \kappa_j^{\rm cell}=-\lambda_j+(d_R-r_j)/2.
\]
Both constants are positive Gaussian normal integrals of the amplitude along
$\mathcal M_j$.  The cell constant uses the full lattice/continuous normal
Hessian, including its mixed blocks.
\end{theorem}
\begin{proof}
A finite tubular cover and a partition of unity reduce the continuous
integral to tangential coordinates $z\in\mathcal M_j$ and normal coordinates
$n$.  Put $n=\eta/\sqrt N$.  The continuous Jacobian is
$N^{-(d_R-r_j)/2}$ and Taylor expansion of the rate gives its positive
normal Gaussian.  Domination permits passage to the limit.  On a pinned
lattice fibre the original $N^{-d_Z/2}$ remains.

For the cell put $\ell=(k-Nk_j)/\sqrt N$.  The lattice mesh has volume
$N^{-d_Z/2}$, so its Gaussian Riemann sum has order $N^{d_Z/2}$.
That factor cancels the single-site lattice prefactor.  The mixed normal
Gaussian must be integrated jointly in $(\ell,\eta)$; conditioning first
would use the appropriate Schur complement.  Collecting the factors in
\eqref{eq:r33-d1-local} proves both expressions.
\end{proof}

For an isolated normalized Gaussian mean, $\lambda_j=d_R/2$ and $r_j=0$.
Thus a summed cell has power zero, while an exact lattice fibre has power
$d_Z/2$.  This verifies both normalizations simultaneously.  Higher-order
coefficients need uniform higher Taylor expansions and matching integrated
remainders; they are not asserted from the leading hypothesis above.

\subsection{Labelled rates and zero prior support}

If a labelled rate $\mathcal I(j,z)$ has independently been proved, define
$\gamma_j=\inf_z\mathcal I(j,z)$ and
$I_j(z)=\mathcal I(j,z)-\gamma_j$.  The contraction is
$\inf_j\mathcal I(j,z)$; after normalization by the common minimum it is
$\min_j(\gamma_j-\min_k\gamma_k+I_j(z))$.
This definition does not use $I_j$ before defining it.
In a finite sum with fixed prior $w$, only
$J_+(w)=\{j:w_j>0\}$ occurs.  Exponentially varying priors require their own
rate shifts and are not covered by the fixed-prior theorem next.

\subsection{Power-drop control of policies near the maximizers}

Let $A$ be a compact metric policy parameter space and let $J_+(w)$ be finite.
Suppose, uniformly in $a\in A$ and $j\in J_+(w)$,
\begin{equation}\label{eq:r33-d1-policy}
 Z_{j,N}(a)=e^{NG_j(a)}N^{-\kappa_j(a)}[c_j(a)+o(1)],
\end{equation}
where $G_j,\kappa_j,c_j$ are continuous and $c_j>0$.
Set $G_* =\max_{j,a}G_j(a)$.  For a phase with maximum $G_*$, write
$A_j^*=\{a:G_j(a)=G_*\}$.  Suppose there are $c,L,p,\alpha>0$ such that
near each nonempty $A_j^*$,
\begin{align}
 G_*-G_j(a)&\ge c\,\operatorname{dist}(a,A_j^*)^p,
 \label{eq:r33-d1-drop}\\
 |\kappa_j(a)-\kappa_j(b)|&\le Ld(a,b)^\alpha.
 \label{eq:r33-d1-holder}
\end{align}
A phase with smaller maximum has a positive uniform exponential gap by
compactness.  Define
\[
 \kappa_* =\min_{j,a\in A_j^*}\kappa_j(a),\qquad
 C_*(a)=\sum_{\substack{j\in J_+(w):\,a\in A_j^*\\
                              \kappa_j(a)=\kappa_*}}w_jc_j(a),
 \quad C_*^{\max}=\max_{a\in A} C_*(a).
\]
The last maximum exists: each active set is closed and its nonnegative
continuous amplitude extended by zero is upper semicontinuous.  At least
one active set is nonempty, so $C_*^{\max}>0$.

\begin{theorem}[Common-policy lexicographic selection with a proved boundary estimate]
\label{thm:r33-d1-selection}
Under the preceding hypotheses,
\[
 \sup_{a\in A}\log\sum_{j\in J_+(w)}w_jZ_{j,N}(a)
 =NG_* -\kappa_*\log N+\log C_*^{\max}+o(1).
\]
No measure is placed on policy space and no policy-volume factor occurs.
\end{theorem}
\begin{proof}
Normalize the sum by $e^{NG_*}N^{-\kappa_*}$.  Since all $\kappa_j$ are
bounded, a phase with deficit $N(G_*-G_j(a))>L_0\log N$ is negligible,
uniformly in $a$, for a sufficiently large fixed $L_0$.
A remaining phase has, by \eqref{eq:r33-d1-drop}, distance at most
$C(\log N/N)^{1/p}$ from $A_j^*$.  Choose a closest $b\in A_j^*$.
Equation \eqref{eq:r33-d1-holder} gives
\[
 \kappa_j(a)\ge\kappa_*-
 C(\log N/N)^{\alpha/p},\qquad
 (\log N)(\log N/N)^{\alpha/p}\longrightarrow0.
\]
Thus a nearby nonmaximizer cannot gain a constant-order advantage by lowering
its polynomial exponent.

For the upper bound, take a maximizing sequence $a_N$ and a convergent
subsequence $a_N\to a$.  Any nonvanishing normalized phase has $a\in A_j^*$.
If $\kappa_j(a)>\kappa_*$ its contribution vanishes; otherwise its limsup
is at most $w_jc_j(a)$ because the exponential deficit is nonnegative and
the possible negative polynomial deficit is $o(1)$ after multiplication by
$\log N$.  There are finitely many phases, so the total limsup is at most
$C_*(a)\le C_*^{\max}$.
For the lower bound fix an $a$ attaining $C_*^{\max}$ in
\eqref{eq:r33-d1-policy}.  Its tied active phases give the matching limit.
Taking logarithms is legitimate since this limit is positive.
\end{proof}

For example, let $A=[-1,1]$, $G_1(a)=-(a-1)^2$,
$G_2(a)=-(a+1)^2$, $c_1=c_2=1$, $\kappa_1=\kappa_2=0$, and
$w_1=w_2=1/2$.  The common-policy amplitude is $1/2$ at either maximizer,
so the constant term is $-\log2$.  Optimizing each phase separately before
summing would give amplitude one.  The leading scalar identity
$\sup_a\max_jG_j(a)=\max_j\sup_aG_j(a)$ remains true, while the constant
order retains the shared-policy restriction.

\subsection{The quantitative price of memory approximation}

Here is a sufficient comparison that applies even to adaptive laws.
Suppose two policies on one chronological experiment have conditional action
densities whose ratio, at every admissible observed past, lies between
$e^{-\epsilon_m}$ and $e^{\epsilon_m}$.  Suppose all other conditional kernels
are identical.  Then their path-density ratio up to time $N$ lies between
$e^{-N\epsilon_m}$ and $e^{N\epsilon_m}$.  Integrating any nonnegative payoff
gives
\[
 |\log Z_N^{\pi}-\log Z_N^{\pi_m}|\le N\epsilon_m.
\]
This follows by multiplying the actual adaptive conditional densities; it is
not an identification with a fixed uncontrolled Feynman--Kac source.

If $\epsilon_m\le Ce^{-\lambda m}$, choosing
$m\ge(1+\eta)\log N/\lambda$ makes this logarithmic error $o(1)$.
But chart constants bounded by $e^{\gamma m}$ then grow as
$N^{\gamma(1+\eta)/\lambda}$, not subpolynomially.  A finite-memory remainder
must absorb this precise power.  Merely summable memory dependence need not
permit both requirements.  Establishing the requisite density-ratio control
and a local expansion for the actual augmented controlled kernel remains
part of the mechanical application.

\subsection{Remaining phase input}

Theorems \ref{thm:r33-d1-laplace} and \ref{thm:r33-d1-selection} are proved
transfer statements with explicit hypotheses.  They do not manufacture the
labelled off-central local density \eqref{eq:r33-d1-local} from an unlabelled
central LLT.  Constructing the physical phases, their normal Hessians, and a
uniform adaptive local theorem is still required for the original theta
phase synthesis.  The present revision retains that target without recording
its missing inputs as established exports.
